
% 10. The Maximization of Fairness and the Convergence of Competition: A Single Symmetric Competitive Track.tex

\section{The Maximization of Fairness and the Collision of Competition: The Symmetric Single Track}
\textbf{The Illusion of Fairness and the Universalization of Competition} \\
We praise modern society as the most advanced era in human history, celebrating the collapse of status and gender barriers as the grand expansion of 'fairness'. We blindly believe that a society stripped of discrimination will bring us a more peaceful and harmonious world.
However, applying the chilling thermodynamic lens of DISTANCE exposes the true physical reality of this noble moral achievement. Fairness is not a flag of peace. It is the most brutal and suffocating process of 'Equalization', shaving away all internal societal barriers (asymmetries) and pouring previously isolated momentums onto a single, identical stage.
This physical truth gives birth to a painful paradox: fairness never eliminates competition; rather, it universalizes it. In past class-based or traditional gender-divided societies, people competed only within their small, isolated domains (local islands), consuming limited energy. But as civilization's massive equalization engine smoothly shaved away the energy gaps of status, institutions, and gender, the momentums of countless isolated individuals were sucked into a single, giant mixer. \\
\textbf{The Ultimate Arena: The SYMMETRIC SINGLE COMPETITIVE TRACK} \\
This explosive fusion of competition is most vividly revealed in the relationship between men and women. For a long time, men and women performed different social roles. This was not merely a result of cultural oppression, but an inevitable functional division of labor designed to efficiently resolve survival momentums in harsh natural environments dependent on muscle power. Men and women lived and were evaluated on completely different tracks.
However, as the Industrial Revolution dismantled the importance of muscle, mass education opened access to knowledge, and information technology erased physical barriers, the walls dividing these two tracks violently collapsed. Furthermore, as robots replaced physical labor and artificial intelligence (AI) began replacing cognitive labor, the physical justifications linking biological sex to economic productivity were completely flattened.
As a result, men and women have been violently thrust onto what DISTANCE defines as the "SYMMETRIC SINGLE COMPETITIVE TRACK". The era of building separate domains while acknowledging differences is over. Modern society has dismantled all asymmetries, completing the most frictionless and perfectly flat ultimate arena in cosmic history, where everyone is judged by a single standard. On this Symmetric Single Competitive Track, men and women now study in the same classrooms, take the same exams, apply for the same jobs, and fiercely collide their momentums within the exact same promotion systems. \\
\textbf{The Collision with Biological Asymmetry} \\
Many praise this state—where all walls of discrimination have crumbled and men and women compete on a perfectly equal footing—as the great victory of civilization. Yet, the lens of DISTANCE warns of the most massive and fatal physical rupture that will inevitably erupt on this stage of perfect equalization.
On this empty arena where social standards have been perfectly symmetrized and everything has become fair, one monumental 'Biological Asymmetry'—which no human political system or economic policy can ever shave down—begins to reveal its terrifyingly massive presence. On this perfectly flattened Symmetric Single Competitive Track, the biological reality of 'pregnancy and childbirth,' which occurs exclusively through the female body, remains an un-erasable physical fact.
How will this innate asymmetry be interpreted by the subjects forced to run at the exact same speed, under the exact same rules, on this flawless Symmetric Single Competitive Track? This is the chilling question that modern civilization must face.